\documentclass[aspectratio=169, 10pt]{beamer}

% --- PAQUETES ---
\usepackage[utf8]{inputenc}
\usepackage[spanish,mexico]{babel}
\usepackage{amsmath, amsthm, amssymb}
\usepackage{hyperref}
\usepackage{graphicx}
\usepackage{float}
\usepackage{animate}
\usepackage{subcaption}
\usepackage{multicol}
\usepackage{verbatim}
\usepackage{color}
\usepackage{ragged2e}
\usepackage[all]{xy}
\usepackage{transparent}
\usepackage{eso-pic}

% --- TEMAS Y ESTILOS PROFESIONALES ---
\usetheme{Madrid}
\usecolortheme{whale}   % Azul profundo elegante en encabezados
\usecolortheme{rose}    % Fondos suavizados y elegantes para bloques
\useinnertheme{rounded} % Esquinas redondeadas y sombras en bloques

% Opciones visuales
\setbeamertemplate{navigation symbols}{} % Oculta los botones de navegación innecesarios
\setbeamercovered{transparent=30}         % Los elementos pausados se ven atenuados en lugar de invisibles

% --- DATOS DE LA PRESENTACIÓN ---
\title[Taller de Herramientas Computacionales]{\Large{\bf Taller de Herramientas Computacionales}}
\subtitle{Grupo 60049 — Semestre 2027-1}

\author[Profesor y Ayudante]{\Large{Miguel Ángel Carrillo Lucía}\and \Large{Karla Elsa Calva Cano}}

\institute[UNAM - FC]{
    \large{Universidad Nacional Autónoma de México}\\
    \large{Facultad de Ciencias}\\
    \large{Licenciatura en Matemáticas Aplicadas}
}

\date{1 de Septiembre de 2026}

% --- FONDO CON ESCUDO (INTACTO) ---
\setbeamertemplate{background}{
  \rule{0pt}{.90\paperheight}%
  \hspace*{.5\paperwidth}%
  \paperwidth1.50mm
  \makebox[1pt][c]{\transparent{0.120} 
\includegraphics[scale=0.56]{logos_unam_ciencias.png}}
}

\begin{document}

% --- DIAPOSITIVA DE TÍTULO ---
\begin{frame}
    \titlepage
\end{frame}

% --- AGENDA ---
\begin{frame}{Agenda del Curso}
    \tableofcontents[sectionstyle=show,subsectionstyle=show/shaded]
\end{frame}

% --- CLASE ---
\section{Información del Curso}
\begin{frame}{Información del Curso}
    \vspace{2mm}
    \begin{block}{Horario, salón de clase y contacto}
    \begin{itemize}\justifying
        \item \textbf{Horario:} Martes y Jueves de 7:00 am a 9:00 am (Grupo 6004). \pause
        \item \textbf{Salón:} 303. Ubicación: Edificio Yelizcalli, Facultad de Ciencias, UNAM. \pause
        \item \textbf{Asesorías remotas:} Programadas a convenir por plataforma o correo electrónico. \pause
        \item \textbf{Asesorías presenciales:} Instituto de Ingeniería (Profesor titular) o Facultad de Ciencias (Profesor ayudante). \pause
         \item Correo Profesor: miguelcarrillo15@ciencias.unam.mx
         \item Correo Profesora: karla\_elsa\_@ciencias.unam.mx
        \item \textbf{Plataforma:} Google Classroom. \textit{Ingresar obligatoriamente con su correo de la Facultad. Código de acceso: aybctacq}
    \end{itemize}
    \end{block}
\end{frame}

% --- EVALUACIÓN ---
\section{Evaluación}
\begin{frame}{Criterios de Evaluación}
    \begin{block}{Desglose de Calificación}
        \begin{itemize} \pause
            \item \textbf{Proyectos (60\%):} Un proyecto por tema. \pause
            \begin{enumerate} 
                \item \textbf{\LaTeX{} (30\%):} Presentación en parejas o equipos de 3 (Temas asignados aleatoriamente. Fecha tentativa de exposición: por definir). \pause
                \item \textbf{Python (30\%):} Análisis de bases de datos (Fecha tentativa: por definir). \pause
            \end{enumerate}
            \item \textbf{Tareas (30\%)} \pause
            \item \textbf{Quizzes (10\%):} Un examen corto al finalizar cada tema. \pause
        \end{itemize}
    \end{block}

    \begin{alertblock}{Requisitos para Examen Final y Asistencia}
        \begin{itemize} \justifying
            \item \textbf{Asistencia obligatoria:} Más de 7 faltas injustificadas enviarán al alumno a examen final.
            \item \textbf{Derecho a Final:} Se requiere haber entregado al menos el 70\% de las actividades académicas. La calificación final se promediará 50\% curso y 50\% examen final.
        \end{itemize}
    \end{alertblock}
\end{frame}

\begin{frame}{Lineamientos de Entregas}
    \begin{block}{Entrega de Trabajos y Código}
    \begin{itemize} \justifying
        \item Habrá retos semanales y ejercicios al inicio de clase. \pause
        \item \textbf{Calidad:} Se evalúa ortografía, gramática y documentación detallada en código. \pause
        \item \textbf{\LaTeX{}:} Se adjuntan archivos comprimidos (\texttt{.zip}) y el documento compilado desde Overleaf (\url{https://es.overleaf.com/}) (\texttt{.pdf}). \pause
        \item \textbf{Python:} Se adjuntan códigos y repositorios en GitHub (Crear cuenta, \url{https://github.com/}). \pause
        \item \textbf{Formato de archivo:} \texttt{THC\_NOMBRE\_TAREA1.zip} \pause
        \item Entrega de calificaciones individual vía correo al finalizar el curso.
    \end{itemize}
    \end{block}

    \begin{alertblock}{Honestidad Académica}
        Trabajo en equipo no es igual a copiar. En caso de copia se asignará \textbf{calificación de cero}. En reincidencia la calificación final será \textbf{NP}.
    \end{alertblock}
\end{frame}

% --- REGLAS ---
\section{Reglas de Convivencia}
\begin{frame}{Reglas de Convivencia y Ética}
    \begin{block}{Trato Institucional}
        \begin{itemize}
            \item Respeto mutuo indispensable entre profesor, ayudante y estudiantes.
        \end{itemize}
    \end{block}

    \begin{block}{Código de Ética}
        \begin{itemize}
            \item Trabajo individual y en equipo con responsabilidad.
            \item Explicar detalladamente todos los desarrollos (ecuaciones, programas, derivaciones).
            \item Citar formalmente las fuentes consultadas para evitar el plagio.
        \end{itemize}
    \end{block}
\end{frame}

% --- TEMARIO ---
\section{Temario}
\begin{frame}{Contenido Sintético del Curso}
\begin{block}{Introducción a \LaTeX{} y Python}
    \begin{enumerate}
        \item \textbf{Introducción al editor de textos en \LaTeX{}.} \pause
        \item \textbf{Introducción a Python:}
        \begin{itemize}
            \item Uso de librerías científicas: NumPy, Matplotlib, Pandas.
        \end{itemize} \pause
        \item \textbf{Aplicaciones y Experimentación:}
        \begin{itemize}
            \item Mínimos cuadrados. Método de Newton-Raphson. Estadística de terremotos en México.
            \item Ley de Enfriamiento de Newton (modelado matemático de fenómenos físicos).
            \item Errores numéricos (redondeo y truncamiento).
            \item Gráficas en 3D.
        \end{itemize}
    \end{enumerate}
\end{block}
\end{frame}

% --- BIBLIOGRAFÍA ---
\section{Bibliografía}

\begin{frame}[fragile]{Bibliografía Recomendada}
\begin{block}{\LaTeX{}}
    \begin{itemize}
        \item Aranda, E. (2013). \textit{Curso de LaTeX}.
        \item Mori, L. F. (2007). \textit{Writing a thesis with LaTeX}.
        \item Oetiker, T. et al. (1995). \textit{La introducción no-tan-corta a LaTeX}.
        \item Gutiérrez, M. (2011). \textit{Notas de Introducción a LaTeX}.
    \end{itemize}
\end{block}

\begin{block}{Python}
    \begin{itemize}
        \item Guzdial, M., \& Ericson, B. (2016). \textit{Introduction to Computing and Programming in Python}. Pearson.
        \item Langtangen, H. P. (2014). \textit{A Primer on Scientific Programming with Python}. Springer.
        \item Punch, W. F., \& Enbody, R. (2010). \textit{The Practice of Computing Using Python}. Addison-Wesley.
    \end{itemize}
\end{block}
\end{frame}

% --- SERVICIOS Y RECURSOS ---
\section{Recursos UNAM}

\begin{frame}{Primer Ingreso}

\begin{center}

\includegraphics[scale=.41]{../../../../../Pictures/Captura de pantalla 2026-07-27 111944.png} 

\url{https://sites.google.com/ciencias.unam.mx/induccion/home}
\end{center}
\end{frame}

\begin{frame}{Primer Ingreso}

\begin{center}
\includegraphics[scale=0.41]{../../../../../Downloads/Captura de pantalla 2026-07-27 115208_2.png} 

\url{https://sites.google.com/ciencias.unam.mx/estuciencias}
\end{center}
\end{frame}


\begin{frame}{Becas}

\begin{center}
\includegraphics[scale=.41]{../../../../../Downloads/Captura de pantalla 2026-07-27 115744_3.png} 

\url{https://sites.google.com/ciencias.unam.mx/becas/abiertas}
\end{center}
\end{frame}



\begin{frame}{Servicios para el Alumnado}
    \begin{block}{Salud Mental y Física}
    \begin{itemize}
        \item \textbf{Centro Médico Universitario:} Optometría, dermatología, odontología, laboratorios.
        \item \textbf{Atención Psicológica:} Servicio ESPORA y atención psiquiátrica.
        \item \textbf{Deportes:} Gimnasios, alberca olímpica, tiro con arco, natación, esgrima, etc.
    \end{itemize}
    \end{block}

    \begin{block}{Oferta Extracurricular y Cultural}
    \begin{itemize}
        \item \textbf{Idiomas (ENALLT):} Chino, alemán, árabe, inglés, entre otros.
        \item \textbf{Cultura UNAM:} Conciertos sinfónicos, obras de teatro, cine y museos.
    \end{itemize}
    \end{block}
\end{frame}

\begin{frame}{Recursos de la Facultad de Ciencias}
    \begin{block}{Préstamo de Equipo y Licencias}
    \begin{itemize}
        \item \textbf{Software Institucional:} MATLAB, Mathematica, Office 365, GitHub Student Pack.
        \item \textbf{Infraestructura:} Megasala de Cómputo y préstamo de tablets.
    \end{itemize}
    \end{block}

    \begin{block}{Actividades Culturales y Deportivas Internas}
    \begin{itemize}
        \item \textbf{Talleres:} Coro, guitarra, teatro, violín.
        \item \textbf{Deportes FC:} Basquetbol, voleibol, fútbol, tochito y dodgeball.
        \item \textbf{Clubes:} Juegos de mesa, cineclub, ajedrez.
    \end{itemize}
    \end{block}
\end{frame}

\end{document}